\begin{figure}[t]
\begin{leangrammar}[sum=GYQkvbUOd0, dflt=x2x8yBvUWM]{Stlc.Tm}{Terms}
\label{gram:Stlc.Tm}\begin{LeanDisplay}
\LeanCategory{M \in \mathit{Term}}{\leanslot{Stlc.Tm}{Terms, with de Bruijn variables}}\\
\LeanEntry{x_{i}}{\leanslot{Stlc.Tm.var}{a variable}}\\
\LeanEntry{\lambda{:}A.\, M}{\leanslot{Stlc.Tm.lam}{an abstraction}}\\
\LeanEntry{M_{1}\; M_{2}}{\leanslot{Stlc.Tm.app}{an application}}\\
\LeanEntry{()}{\leanslot{Stlc.Tm.unit}{the unit value}}
\end{LeanDisplay}
\end{leangrammar}

\begin{leanrules}[sum=mXx-JpINJ1]{Stlc.IsValue}{Values}
\begin{mathpar}

\inferrule[(Lam)]
  { }
  { (\lambda{:}A.\, M) \in \mathit{Val} }

\and

\inferrule[(Unit)]
  { }
  { () \in \mathit{Val} }

\end{mathpar}
\end{leanrules}

\begin{leanrules}[sum=5cW2xFjnlQ]{Stlc.Step}{Reduction}
\begin{mathpar}

\inferrule[(Beta)]
  { }
  { (\lambda{:}A.\, M)\; V \longrightarrow M[V] }

\and

\inferrule[(App-Left)]
  {
    M_{1} \longrightarrow M_{1}'
  }
  { M_{1}\; M_{2} \longrightarrow M_{1}'\; M_{2} }

\and

\inferrule[(App-Right)]
  {
    M_{2} \longrightarrow M_{2}'
  }
  { V\; M_{2} \longrightarrow V\; M_{2}' }

\end{mathpar}
\end{leanrules}

\begin{leantheorem}[sum=8jx5_cPcFh]{Values are final}{Values are final}
\label{thm:Values-are-final}
\LeanProse{Let $V \in \mathit{Val}$ and $M \in \mathit{Term}$. Then $\neg (V \longrightarrow M)$.}
\end{leantheorem}

\begin{leansource}[sum=zIy9s_XPen]{Stlc.values_final}{The source}
\LeanCode{%
\LeanKeyword{theorem} \LeanConst{values\_final} \{\LeanBound{v} \LeanBound{e} : \LeanConst{Tm}\} (\LeanBound{hv} : \LeanConst{IsValue} \LeanBound{v}) : ¬\LeanConst{Step} \LeanBound{v} \LeanBound{e}}
\end{leansource}
\caption{The running example: a simply typed $\lambda$-calculus, as the
paper's reader sees it, and the source of its theorem. All five blocks are
generated.}
\label{fig:example}
\end{figure}
